% Proposed prose blocks for paper/main.tex. Not included automatically.
% A: replace the abstract body, retaining the existing keyword block.
We study end-of-day execution through a continuous limit-order book followed
by a closing auction. The trader begins with positive inventory and sells
during continuous trading; signed auction schedules subsequently allow
inventory adjustment at clearing. We develop a reinforcement-learning framework
that combines a projected clearing-price signal with intermediate feedback
for auction decisions whose cash consequences are delayed until settlement.
We compare a baseline DQN with projected continuous-proposal DDPG, TD3 and SAC
policies. Headline policies train on a shaped objective, while selection and
evaluation use marked-to-market PnL net of fees and terminal inventory risk.
Experiments with synthetic rough-Heston prices and historical midprice inputs
inside the same simulated market show improved risk-adjusted performance
relative to stylized AS and TWAP references. The continuous methods obtain
larger realized auction contributions than DQN, with historical-setting gains
arising mainly from terminal inventory-risk relief. Matched ablations show
that dense auction credit improves economic learning. The clearing forecast
is predictively informative, although its incremental economic value varies
by learner. Additional forecast-relative reward preferences have no established
average benefit once learning guidance is held fixed. These findings identify
both the usefulness and the limits of auction-informed learning for liquidation.

% B: introduction question/contribution paragraphs; see the integration plan.
The central question is whether a learning agent can use information about a
future closing auction to balance continuous execution with end-of-day inventory
management. This requires connecting two decisions: how much inventory to sell
during continuous trading and how to adjust the residual position through an
auction whose cash consequences are realized only at clearing. Our framework
provides a projected clearing-price signal, an explicit auction control problem
and intermediate learning guidance. We study DQN as a baseline and DDPG, TD3
and SAC as continuous-proposal extensions projected onto the same executable
action grid.

The numerical study addresses three questions. First, is the projected clearing
signal informative, and does supplying it improve economic decisions? Second,
does intermediate training guidance help policies learn from delayed auction
outcomes? Third, does access to the closing auction improve liquidation and
terminal inventory management? We answer these questions using matched
experiments alongside comparisons with stylized reference policies. Throughout,
we distinguish training feedback from economic cash flows and evaluate all
policies using the same risk-adjusted economic criterion.

% C: add after the definition/discussion of the unshaped economic criterion.
\paragraph{The role of training guidance.}
The fictive rewards are learning signals rather than exchange cash flows.
Their purpose is to guide actions before auction settlement, and their usefulness
must therefore be assessed through the economic behavior of the resulting policy.
Excluding them from $\bar J_\lambda$ prevents an artificial increase in the
reported score; it does not prevent a better learned policy from improving that
score. We distinguish the additional preferences defining the shaped objective
$J$ from potential-based temporal credit used in replay. The latter redistributes
feedback across decision times with a terminal correction that preserves the
episode objective up to a policy-independent constant. Our ablations examine
the complete training scheme, temporal credit assignment, and the additional
preferences separately.

% D: empirical protocol and first results subsection, before the bibliography.
\subsection{Evaluation Protocol and Economic Performance}
\label{subsec:empirical_protocol}
The headline study uses four learning algorithms and ten independent training
seeds in a synthetic setting and five settings driven by historical midprices.
Order flow, auction updates and clearing remain simulated in all settings.
Policies are selected on separate economic validation paths and evaluated on
100 additional test episodes per selected policy. We average outcomes within
each training seed and give seeds equal weight. Matched differences pair the
evaluation environments before averaging; pointwise 95\% bootstrap intervals
resample training-seed means. The initial 440-run campaign was followed by two
declared 40-run synthetic comparisons with literal economic cash-flow training
and with conditioned economic training at the headline auction anchor.
The full specifications and the chronology of these follow-ups are reported
in the experimental appendix.

Headline policies outperform both stylized reference policies on mean
$\bar J_\lambda$ in all six settings, with positive pointwise paired intervals.
The advantage should be interpreted through both execution PnL and inventory
risk. The implemented risk-neutral AS reference often earns greater raw PnL
while retaining more terminal risk. The learned policies improve the risk-adjusted
trade-off rather than universally maximizing raw PnL. Likewise, the continuous
methods are competitive with DQN on total performance, but the results do not
establish a universal algorithm ranking. Negative economic means on the MSFT
midprice setting are retained in the reported results.

\subsection{Auction Access, Execution Value and Terminal Inventory}
\label{subsec:auction_value}
The auction is a final execution opportunity, and its value should be measured
against an explicit economic counterfactual. For a fixed realized CLOB trajectory,
the contribution of the subsequent learned auction policy relative to submitting
no auction orders is
\[
Z_{\tau^\cl}\bigl(S_{\tau^\cl}^\cl-S_{\tau^\op}^{\mathrm{mid}}\bigr)
-\sum_{t_i\geq\tau^\op}d_{t_i}c_{t_i}
+\lambda\bigl(I_{\tau^\op}^{2}-I_{\tau^\cl}^{2}\bigr).
\]
This separates execution price edge, cancellation fees and inventory-risk
relief. The direct contribution is positive on average for every headline
market--method combination; its pointwise interval excludes zero in 23 of
the 24 combinations, with synthetic DQN the exception. In the historical
settings, equal-ticker mean contributions are approximately 0.39, 1.48,
1.35 and 1.54 basis points for DQN, DDPG, TD3 and SAC. Most of this value
comes from reduced terminal inventory risk. The continuous methods obtain
larger realized contributions than DQN, including the effect of their own
endogenous opening inventories.

A separate experiment retrains policies without auction access, terminating
at the CLOB boundary with residual marking and the inventory penalty. Relative
to this control, auction access improves the economic objective by 0.22,
0.56, 0.48 and 0.61 basis points for DQN, DDPG, TD3 and SAC, respectively.
The continuous-method intervals are positive; the DQN interval includes zero.
This comparison holds economic training and explicit forecast access fixed.
Together, the retrained access experiment and the direct decomposition show
that the learned continuous policies use the auction to economic advantage.
Because the auction permits signed positions, this need not imply lower
absolute inventory on every path or complete liquidation.

% E: forecast results subsection.
\subsection{Forecast Information and Economic Decisions}
\label{subsec:forecast_value}
The projected clearing signal has lower mean absolute and root mean squared
prediction errors than contemporaneous midprice in both phases across all
six settings. These errors use eventual simulated auction clearing as the
target. For example, synthetic mean absolute error is 0.1800 versus 0.1925
during continuous trading and 0.0488 versus 0.1403 during the auction.
This validates predictive information within the modeled market; it does
not constitute an auction-price forecasting test on observed exchange data.

To examine decision value, we compare access to the explicit $H_t^\cl$
observation at a common fixed quote anchor, economic objective and learning
guidance. DDPG gains 0.72 basis points, with interval $[0.12,1.36]$.
The corresponding DQN, TD3 and SAC effects are 0.17, 0.15 and $-0.17$
basis points, with intervals including zero. Forecast information is therefore
predictively useful, while its incremental economic value depends on the
learner and the information already supplied. The H-off policies retain other
clearing-relevant state variables. A separate anchoring comparison shows that
centering a bounded action grid on the forecast need not improve performance;
its complete results are reported in the appendix. Prediction quality, access
to a prediction and the geometry of executable quotes are distinct properties.

% F: training-guidance results subsection.
\subsection{Learning from Delayed Auction Outcomes}
\label{subsec:auction_learning}
During the auction, submissions and cancellations affect the eventual position
before their cash consequences are realized. Dense feedback addresses the
resulting credit-assignment difficulty. In the matched auction-potential
comparison, dense credit improves final economic performance by 1.18, 1.69,
2.30 and 0.85 basis points for DQN, DDPG, TD3 and SAC, respectively, with
positive pointwise intervals for every algorithm. Economic validation curves
also show positive paired differences at episode 250. Both arms retain common
CLOB conditioning and optimize the economic objective, so this experiment
isolates an auction-credit effect rather than an extra economic cash flow.

The complete headline training scheme, including weighted fictive rewards,
also outperforms literal economic cash-flow training for all four methods
and all ten paired seeds. However, that comparison changes the entire reward
construction. To identify the additional manuscript preferences, the final
comparison retains the headline's H-centred quotes and all reward conditioning
while training on $\bar J_\lambda$. The J-minus-economic effects are
$-0.06$, $-0.06$, $0.19$ and $-0.20$ basis points, respectively; all intervals
include zero. Conditioned economic policies improve selected validation over
initialization in 39 of 40 runs and obtain positive mean auction contributions.
The evidence supports intermediate guidance for delayed auction decisions,
while additional forecast-relative preferences have no established marginal
economic benefit at the tested calibration. This does not prove equivalence
of the objectives or their learned policies.

% G: concluding section after the empirical results.
\section{Conclusion}
An explicit closing auction gives a learning agent an additional means to
manage an initial inventory position. Our framework connects continuous
execution, a projected clearing signal and delayed auction settlement, and
produces policies that improve risk-adjusted performance relative to the stated
stylized references. The continuous-proposal methods obtain greater realized
auction contributions than the DQN baseline, primarily through terminal
inventory-risk relief in the historical-midprice settings. Matched experiments
show that dense auction credit improves economically evaluated learning.
The forecast is informative, although its incremental policy value varies
across methods, and the additional manuscript preferences do not have a
demonstrated average advantage once conditioning is shared. These findings
support auction-informed learning for end-of-day execution while delimiting
the claims justified by a sequential simulated market and finite training data.
